\documentclass[10pt,a4paper]{amsart}
\usepackage[margin=19mm]{geometry}
\usepackage{amsmath,amssymb,mathtools}
\usepackage[scr=rsfs,cal=euler]{mathalfa}
\usepackage{xfrac}
\usepackage[unicode,hidelinks,bookmarks=false]{hyperref}
\newcommand{\ie}{i.e.\ }
\newcommand{\cf}{cf.\ }
\newcommand{\resp}{respectively\ }
\newcommand{\Subsection}{\subsection}
\providecommand{\nobreakdash}{\nobreak\hbox{-}}
\setlength{\emergencystretch}{2em}
\multlinegap0pt
\usepackage{bm}
\usepackage{bbm}
\usepackage{dsfont}
\usepackage{xspace}
\usepackage{cancel}
\usepackage{mathtools}
\usepackage{amsthm}
\makeatletter
\@ifundefined{theorem}{\newtheorem{theorem}{Theorem}}{}
\makeatother
\title[De Bruijn Tori Without Zeros: A Field-Theoretic Perspective]{De Bruijn Tori Without Zeros: A Field-Theoretic Perspective}
\author{Ming Hsuan Kang, Yu Hsuan Hsieh}
\date{Source: arXiv:2506.19605v1 (CC BY 4.0)}
\begin{document}
\maketitle

\noindent\emph{Source and licence.} De Bruijn Tori Without Zeros: A Field-Theoretic Perspective. Version arXiv:2506.19605v1 (CC BY 4.0), distributed under the Creative Commons Attribution 4.0 International licence, as recorded in the arXiv licence metadata for this exact version and shown on the abstract page https://arxiv.org/abs/2506.19605v1. Full rights notice: licenses/arxiv-2506.19605-CC-BY-4.0.txt.\par\smallskip

\noindent\textit{Editorial scope.} Counted unit: the Theorem whose source label is \texttt{None} in the article (source0.tex lines 190--192, source label \texttt{None}), delivered together with the article's own complete original proof of it (source0.tex lines 193--220, one \texttt{proof} environment of 28 lines). No mathematics is written, repaired or re-derived: statement and proof are the article's own text line for line. Required context retained uncounted: none. Every definition, display and label of those lines is retained unchanged. Omitted from this package and not redistributed: 1--189, 221--777. Nothing inside a retained excerpt is omitted or altered: every retained body is delivered exactly as the article's lines are written, so no omission receipt of any kind accompanies this package. Citations: the retained excerpts carry no citation command at all, so no bibliography entry of the article is transcribed and no reference is redistributed. Reference adaptation: none was needed and none was made; every reference inside the delivered unit resolves inside it, so no unresolved reference or citation is produced. Printed slips kept literal: no printed word, symbol, hypothesis or number of the counted statement or of its proof was altered. Layout and class: the article's own class is \texttt{article} with packages including graphicx, amssymb, amsmath, mathtools, amsthm, amsfonts, tikz-cd, mathrsfs; the wrapper is the lane's pinned \texttt{amsart} editorial wrapper at 10pt on A4, which restates the theorem-like environments the retained text needs and adds the article's own macro definitions that the retained text needs (none), with extra wrapper packages (bm, bbm, dsfont, xspace, cancel, mathtools). The delivered numbering is the unit's own. Assets: no retained span contains a figure environment, a graphics-inclusion command, a PDF-page inclusion, a table or a TikZ picture; no image file is copied into this package, the package declares no asset dependency and no image file is redistributed with it. Licence: version arXiv:2506.19605v1 of this article is distributed under the Creative Commons Attribution 4.0 International licence, as recorded in the arXiv licence metadata for this exact version and shown on the abstract page https://arxiv.org/abs/2506.19605v1. A full rights notice is delivered at licenses/arxiv-2506.19605-CC-BY-4.0.txt. Characters: every retained excerpt is pure ASCII, so the delivered engine, XeLaTeX with its default Latin Modern text font, reports no missing character.
\begin{theorem}
Let \( B = (B_{i,j}) \) be the De Bruijn torus defined by \( B_{i,j} = \psi(\beta^i \gamma^j) \), where \( \psi: \mathbb{F}_{p^n} \to \mathbb{F}_p \) is a fixed nonzero \(\mathbb{F}_p\)-linear map. Then a pattern \( S\) of size $n$ is a sampling pattern if and only if  it is a basis pattern.
\end{theorem}

\begin{proof}
Define an \(\mathbb{F}_p\)-linear map
\[
\Phi: \mathbb{F}_{p^n} \to \mathbb{F}_p^{S}, \quad
\Phi(z) := \left( \psi(x z) \right)_{x \in A|_S}.
\]
We will show that \(\Phi\) is an isomorphism if and only if \( A|_S \) forms a basis of \( \mathbb{F}_{p^n} \) over \(\mathbb{F}_p\).

\textbf{(\(\Rightarrow\))} Suppose \(\Phi\) is a linear isomorphism. Then the coordinate functionals
\[
z \mapsto \psi(x z), \quad \text{for each } x \in A|_S,
\]
must form an \(\mathbb{F}_p\)-linearly independent set in the dual space \( \mathrm{Hom}_{\mathbb{F}_p}(\mathbb{F}_{p^n}, \mathbb{F}_p) \). Suppose for contradiction that the set \( A|_S \) is linearly dependent, so there exists a nontrivial relation
\[
\sum_{x \in A|_S} a_x x = 0, \quad \text{with } a_x \in \mathbb{F}_p, \text{ not all zero}.
\]
Then for every \( z \in \mathbb{F}_{p^n} \), we have
\[
\sum_{x} a_x \psi(x z) = \psi\left( \left( \sum_{x} a_x x \right) z \right) = \psi(0) = 0,
\]
which implies that the linear combination \( \sum_{x} a_x \psi(x z) \) vanishes identically as a function of \( z \). This contradicts the assumption that the coordinate functionals are linearly independent. Hence, \( A|_S \) must be linearly independent. Since \( |S| = n = \dim_{\mathbb{F}_p} \mathbb{F}_{p^n} \), it follows that \( A|_S \) is a basis.

\textbf{(\(\Leftarrow\))} Conversely, suppose \( A|_S \) is an \(\mathbb{F}_p\)-basis of \( \mathbb{F}_{p^n} \). We prove that \(\Phi\) is an isomorphism by showing it is injective.

Assume \( \Phi(z) = 0 \) for some \( z \in \mathbb{F}_{p^n} \). Then \( \psi(x z) = 0 \) for all \( x \in A|_S \). Since \( A|_S \) is a basis and \( z \ne 0 \), the set \( \{ x z \mid x \in A|_S \} \) also forms a basis of \( \mathbb{F}_{p^n} \) (as multiplication by a nonzero scalar is an automorphism of the vector space). Therefore, \( \psi \) vanishes on a basis and hence on all of \( \mathbb{F}_{p^n} \), which contradicts the assumption that \( \psi \) is nonzero. Thus \( z = 0 \), and \(\Phi\) is injective.

Since both the domain and codomain of \(\Phi\) are \( \mathbb{F}_p \)-vector spaces of dimension \( n \), injectivity implies that \(\Phi\) is an isomorphism.
\end{proof}

\end{document}
